\chapter{Regression and Classification}

Linear and logistic regression are the simplest ML models, but they
contain — in miniature — every mathematical idea used by deeper models:
a parametric function, a loss function, and gradient-based fitting.

\section{Linear Regression}

The model predicts a continuous target as a linear combination of features:
\[
  \hat y = \mathbf{w}^\top \mathbf{x} + b.
\]

\begin{definition}[Mean Squared Error]
\[
  J(\mathbf{w}, b) = \frac{1}{n}\sum_{i=1}^n \left(\hat y_i - y_i\right)^2
\]
\end{definition}

Because $J$ is convex and quadratic in $\mathbf{w}$, it has a closed-form
solution — the \textbf{normal equation}:
\[
  \mathbf{w}^\star = (X^\top X)^{-1} X^\top \mathbf{y}
\]
In practice, gradient descent is often used instead, especially when $X^\top X$
is large or ill-conditioned to invert directly.

\figw{linear_regression_fit.png}{0.68}{Least-squares fit to noisy data. Gray lines show residuals — the vertical distances MSE penalizes.}

\section{Logistic Regression}

For binary classification, we squash the linear output through a sigmoid
to get a probability:
\[
  \hat y = \sigma(\mathbf{w}^\top\mathbf{x} + b) = \frac{1}{1+e^{-(\mathbf{w}^\top\mathbf{x}+b)}} \in (0,1)
\]

\begin{definition}[Binary Cross-Entropy Loss]
\[
  J(\mathbf{w}, b) = -\frac{1}{n}\sum_{i=1}^n \Big[y_i \log \hat y_i + (1-y_i)\log(1-\hat y_i)\Big]
\]
\end{definition}

Unlike linear regression, logistic regression has no closed-form solution
— $J$ is convex, but nonlinear, so it's fit with gradient descent. The
gradient turns out to have an elegantly simple form:
\[
  \nabla_{\mathbf{w}} J = \frac{1}{n} X^\top (\hat{\mathbf{y}} - \mathbf{y})
\]
which looks just like the linear regression gradient, with $\hat y$ now
being the sigmoid output instead of the raw linear prediction.

\figw{logistic_decision_boundary.png}{0.62}{Logistic regression fit to two Gaussian-distributed classes. Color shows predicted probability $P(y=1\mid x)$; the black curve is the decision boundary at $P=0.5$.}

\section{Regularization}

To prevent overfitting, a penalty on the weights is added to the loss:
\[
  J_{\text{reg}}(\mathbf{w}) = J(\mathbf{w}) + \lambda \underbrace{\|\mathbf{w}\|_2^2}_{\text{L2 (Ridge)}}
  \qquad\text{or}\qquad
  J_{\text{reg}}(\mathbf{w}) = J(\mathbf{w}) + \lambda \underbrace{\|\mathbf{w}\|_1}_{\text{L1 (Lasso)}}
\]

\begin{keyidea}
L2 regularization shrinks weights smoothly toward zero (but rarely exactly
zero), while L1 regularization tends to push many weights to \emph{exactly}
zero — producing sparse models that effectively perform feature selection.
This difference comes directly from the geometry of the $\ell_1$ vs.\
$\ell_2$ norm balls.
\end{keyidea}
